\scanpage{49}
$U$ sss $x\in U^{\unsure{\nu|v}}$, \deleted{\uline{donc}} cela implique
$x\in U_i^\nu$ pour $i$ convenable i.e. les $x_\alpha\in U_i$.
O.K.

On est ainsi ramené au cas où $X$
est \uline{affine} sur $Y$. Mais alors on est
\unsure{=} ramené au cas où $X=\operatorname{Spec}(\deleted{\unsure{S}}\unsure{\mathrm V F|F|\underline F})$, où
$F$ est un faisceau loc. libre de rang \unsure{$n$}
sur $T$. On sait qu'alors
\[
f_\ast(X/T)=\operatorname{Spec}(f_\ast \unsure{\mathrm V F|F|\underline F}).
\]

\deleted{\uline{\unsure{Corollaire|Contre-exemple}} $\unsure{f_\ast}$ \unsure{Si}}

\uline{Remarque} Il n'est pas vrai en général que
la dim. des fibres de $f_{\deleted{\unsure{\ast}}\ast}(X/T)/S$
soit majorée par $\nu$ fois la dim
des fibres de \deleted{$\unsure{f_\ast}$} $X/T$. Ex:
$S=\operatorname{Spec}(k)$, $T$\deleted{$=\operatorname{Spec}\unsure{k[t]/t^2|k[t]/t^p|k[E]/t^2|k[E]/t^p}$} est purement
radiciel sur $S$, $X=T\times_S T$ i.e. $f_\ast(X/T)=$
$\underline{\operatorname{Hom}}_S(T,T)$. Ce dernier peut être de
\uline{dim} \inserted{$>0$ \deleted{$\unsure{0}$}} sur $k$. Exemple: le groupe des
automorphismes de la $k$-algèbre \deleted{$\unsure{A|L}$} est
de dim $>0$.
% slanted
\margin{\unsure{voir pour ceci|voir précisé} plus bas}
% archive sheet number: 48
